% ============================================================================
%  Chapitre 4, partie II : Algorithme de simulation
% ============================================================================

% ----------------------------------------------------------------------------
\section{Algorithme de simulation - Orientation en angles d'Euler}
% ----------------------------------------------------------------------------
On peut se demander pourquoi on effectue une simulation avec les angles
d'Euler pour décrire l'orientation, alors que l'on abandonne cette
piste juste après : la simulation ne tient que tant que
$|\cos\theta| > \varepsilon$. La raison est simple : c'est la
traduction directe de tout ce que l'on a développé auparavant. Cette
simulation montre ce que donne la dynamique : la trajectoire de $G$,
les angles, les réactions. L'interprétation des quaternions, en
revanche, n'est pas directe, et intuitive et cette simulation permet donc de fixer
ce que l'on attend du simulateur en quaternions. C'est pourquoi on ne
fait pas, dans cette partie, d'analyse de convergence détaillée.
% ----------------------------------------------------------------------------
\subsection{Système global}

On note $\mat{K}_\alpha = \mat{M}_\alpha + m\,\mat{A}_\alpha^T\mat{A}_\alpha$.
D'après le corollaire 4.6, l'état $(r,\alpha,\dot\alpha)$ vérifie le
système du premier ordre
\[
    \ddt\begin{pmatrix} r\\ \alpha\\ \dot\alpha \end{pmatrix}
    = \begin{pmatrix}
        -\,\mat{A}_\alpha\,\dot\alpha\\
        \dot\alpha\\
        \mat{K}_\alpha^{-1}\bigl(-\,\vect{h}_\alpha
        - m\,\mat{A}_\alpha^T\dot{\mat{A}}_\alpha\,\dot\alpha\bigr)
    \end{pmatrix}
\]
Les deux dernières lignes ne dépendent pas de $r$. On résout donc
d'abord $(\alpha,\dot\alpha)$ seul, puis on en déduit $r$ et
$\vect{\lambda}$.

% ----------------------------------------------------------------------------
\subsection{Sous-système en orientation}

\begin{proposition}[Sous-système en orientation]
    On pose $Y = (\alpha,\dot\alpha)$, qui vérifie
\[
    \dot Y = f(Y) = \begin{pmatrix} \dot\alpha\\ \ddot\alpha \end{pmatrix},
    \qquad
    \mat{K}_\alpha\,\ddot\alpha = -\,\vect{h}_\alpha
    - m\,\mat{A}_\alpha^T\dot{\mat{A}}_\alpha\,\dot\alpha
\]
\end{proposition}
On peut résoudre numériquement ce système avec les conditions initiales de $\alpha$ et $\dot{\alpha}$, ce qui permet de calculer $\ddot{\alpha}$ par résolution d'un système d'équations. Ensuite, une fois obtenu ces trois informations, on peut calculer la position de centre de masse et les forces de réaction.

% ----------------------------------------------------------------------------
\subsection{Sous-système en position et réactions}

\begin{proposition}[Sous-système en position et réactions]
    Une fois $(\alpha,\dot\alpha,\ddot\alpha)$ connu, le reste se calcule facilement : 
\[
    \dot r = -\,\mat{A}_\alpha\,\dot\alpha,
    \qquad
    \vect{\lambda} = -\,m\,\bigl(\mat{A}_\alpha\,\ddot\alpha
    + \dot{\mat{A}}_\alpha\,\dot\alpha\bigr)
\]
\end{proposition}


% ----------------------------------------------------------------------------
\subsection{Schéma de Runge-Kutta d'ordre 4 (RK4) pour la résolution de la dynamique}

\subsubsection{Schéma RK4}

Pour $\dot Y = f(Y)$, avec un pas $h$ et $Y_k \approx Y(t_k)$,
$t_k = k\,h$,
\[
    \begin{aligned}
        \vect{p}_1 &= f(Y_k), &
        \vect{p}_2 &= f\bigl(Y_k + \tfrac{h}{2}\,\vect{p}_1\bigr),\\
        \vect{p}_3 &= f\bigl(Y_k + \tfrac{h}{2}\,\vect{p}_2\bigr), &
        \vect{p}_4 &= f\bigl(Y_k + h\,\vect{p}_3\bigr)
    \end{aligned}
    \qquad
    Y_{k+1} = Y_k + \tfrac{h}{6}\,\bigl(\vect{p}_1 + 2\,\vect{p}_2
    + 2\,\vect{p}_3 + \vect{p}_4\bigr)
\]
L'erreur est en $O(h^5)$ sur un pas et en $O(h^4)$ sur un intervalle
fixé.

\subsubsection{Schéma RK4 pour simuler le système}

On se donne la condition initiale $(r_0,\alpha_0,\dot\alpha_0)$, avec
$\cos\theta_0 \neq 0$.

\paragraph{Étape 1 : orientation.} Avec $Y_k = (\alpha_k,\dot\alpha_k)$,
on évalue $f$ en quatre états intermédiaires du pas,
\[
    Y^{(1)} = Y_k,
    \qquad
    Y^{(2)} = Y_k + \tfrac{h}{2}\,\vect{p}_1,
    \qquad
    Y^{(3)} = Y_k + \tfrac{h}{2}\,\vect{p}_2,
    \qquad
    Y^{(4)} = Y_k + h\,\vect{p}_3
\]
où $Y^{(i)} = (\alpha^{(i)},\dot\alpha^{(i)})$ et
\[
    \vect{p}_i = f\bigl(Y^{(i)}\bigr)
    = \begin{pmatrix} \dot\alpha^{(i)}\\ \ddot\alpha^{(i)} \end{pmatrix},
    \qquad
    \mat{K}_\alpha\bigl(\alpha^{(i)}\bigr)\,\ddot\alpha^{(i)}
    = -\,\vect{h}_\alpha\bigl(\alpha^{(i)},\dot\alpha^{(i)}\bigr)
    - m\,\mat{A}_\alpha^T\dot{\mat{A}}_\alpha\,\dot\alpha\,\big|^{(i)}
\]
On obtient alors le nouvel état $Y_{k+1} = (\alpha_{k+1},\dot\alpha_{k+1})$ :
\[
    \alpha_{k+1} = \alpha_k + \tfrac{h}{6}\,\bigl(\dot\alpha^{(1)}
    + 2\,\dot\alpha^{(2)} + 2\,\dot\alpha^{(3)} + \dot\alpha^{(4)}\bigr),
    \qquad
    \dot\alpha_{k+1} = \dot\alpha_k + \tfrac{h}{6}\,\bigl(\ddot\alpha^{(1)}
    + 2\,\ddot\alpha^{(2)} + 2\,\ddot\alpha^{(3)} + \ddot\alpha^{(4)}\bigr)
\]

\paragraph{Étape 2 : accélération.} On résout le système réduit au
nouvel état :
\[
    \mat{K}_\alpha(\alpha_{k+1})\,\ddot\alpha_{k+1}
    = -\,\vect{h}_\alpha(\alpha_{k+1},\dot\alpha_{k+1})
    - m\,\mat{A}_\alpha^T\dot{\mat{A}}_\alpha\,\dot\alpha\,\big|_{k+1}
\]
C'est la seconde moitié de $f(Y_{k+1})$, réutilisée comme $\vect{p}_1$
au pas suivant.

\paragraph{Étape 3 : position et réactions.} En parallèle avec étape 1, on calcule des poids RK4 pour $r$ avec 
$\vect{g}_i = -\,\mat{A}_\alpha\bigl(\alpha^{(i)}\bigr)\,
\dot\alpha^{(i)}$ tels que : 
\[
    r_{k+1} = r_k + \tfrac{h}{6}\,\bigl(\vect{g}_1 + 2\,\vect{g}_2
    + 2\,\vect{g}_3 + \vect{g}_4\bigr),
    \qquad
    \vect{\lambda}_{k+1} = -\,m\,\bigl(\mat{A}_\alpha(\alpha_{k+1})\,
    \ddot\alpha_{k+1}
    + \dot{\mat{A}}_\alpha\,\dot\alpha\,\big|_{k+1}\bigr)
\]


% ----------------------------------------------------------------------------
\subsection{Conditions initiales et choix du pas}

\subsubsection{Conditions initiales}
\begin{proposition}[Conditions initiales]
    Les conditions initiales de cette dynamique sont $(r_0, \alpha_0, \dot{\alpha}_0)$.
\end{proposition}

On choisit librement la position $r_0$, l'orientation $\alpha_0$ avec
$\cos\theta_0 \neq 0$, et les vitesses angulaires $\dot\alpha_0$. En
pratique, on se donne plutôt la vitesse angulaire $\vect{\Omega}_0$ dans
$\rep{s}$ et
\[
    \dot\alpha_0 = \mat{E}_s(\alpha_0)^{-1}\,\vect{\Omega}_0,
    \qquad
    \det\mat{E}_s = \cos\theta_0 \neq 0
\]
\begin{remarque}[Sur les vitesses angulaires]
    On peut aussi se donner la vitesse angulaire $\vect{\omega}_0$
    dans le repère corps. Cependant, on préfère le repère sol, car ses axes sont fixes et visuels et chaque composante de $\vect{\Omega}_0$ a un sens direct :
    $\Omega_x$ et $\Omega_y$ fixent le roulement initial,
    $\dot r_0 = a\,(\Omega_y, -\Omega_x)$, et $\Omega_z$ la rotation
    autour de la verticale. Les axes du corps, selon nos conventions, sont les axes principaux d'inertie :
    ils dépendent de la répartition de masse à l'intérieur de la
    sphère, et il est difficile de les repérer visuellement.
\end{remarque}

Par contre, la vitesse de $G$ n'est pas libre, le roulement impose
$\dot r_0 = -\,\mat{A}_\alpha(\alpha_0)\,\dot\alpha_0$. Enfin,
$\ddot\alpha_0$ et $\vect{\lambda}_0$ s'obtiennent directement en faisant les mêmes calculs que l'étape 2.

\subsubsection{Choix du pas}

\textbf{A discuter !!!!!}
Il faut noter qu'avec la modélisation en angles d'Euler, les conditionnements de notre système s'explose lorsqu'on a gimbal lock.

% ----------------------------------------------------------------------------
\subsection{Contrôle de la simulation}

Le schéma ne garantit pas la physique : on contrôle après coup que la
trajectoire calculée respecte ce que le modèle impose.

\subsubsection{Quantités conservées}

Ni le poids ni la force de réaction travaille, l'énergie cinétique doit rester donc constante.
\[
    T = \tfrac12\,m\,\dot r^T\dot r
    + \tfrac12\,\vect{\omega}^T\mat{I}_b\,\vect{\omega}
\]
est conservée. 

De plus, la réaction s'applique en $P$ et le poids passe
par $P$, leur moment en $P$ est nul, donc le moment cinétique au point
de contact 
\[
    \vect{L}_P = \mat{I}_P\,\vect{\Omega},
    \qquad
    \mat{I}_P = \mat{I}_G + m\,a^2\bigl(\Id - \vect{e}_z\vect{e}_z^T\bigr) = \mat{R} \mat{I}_b \mat{R} ^T + m\,a^2\bigl(\Id - \vect{e}_z\vect{e}_z^T\bigr)
\]
doit être conservé dans $\rep{s}$.

\begin{remarque}[Racine de l'anomalie de trajectoire de la sphère de Chaplygin]
    On a ici que $\vect{L}_P = \mat{I}_P \vect{\Omega}$ est conservée donc $\vect{\Omega} = \mat{I}_P ^{-1} \vect{cte}$ ou $\vect{\Omega} = {(\mat{R} \mat{I}_b \mat{R} ^T + m\,a^2\bigl(\Id - \vect{e}_z\vect{e}_z^T\bigr) )}^{-1}\vect{cte}$. Pour une sphère isotrope, $\mat{I}_P$ ne dépend pas de $\mat{R}$, donc $\vect{\Omega}$ reste une constante au long de la trajectoire, alors que si la sphère est anisotrope, $\mat{I}_P$ dépend de $\mat{R}(t)$, donc $\vect{\Omega}$ varie au long de la trajectoire. 

    D'un autre point de vue, si la sphère est isotrope ($\mat{I}_b = i\,\Id$),
    le terme gyroscopique $\vect{\omega}\wedge(\mat{I}_b\vect{\omega})$ est nul,
    $\mat{I}_P$ ne dépend pas de l'orientation et $\vect{\Omega}$ reste constant,
    la sphère roule ainsi en ligne droite. Sinon, le terme gyroscopique est non nul :
    $\vect{L}_P$ reste fixe dans $\rep{s}$ alors que l'inertie tourne avec le
    corps, donc $\vect{\omega}$ se réoriente dans $\rep{b}$, et $\vect{\Omega}$,
    puis la direction de $\dot{r}$.
\end{remarque}
On suit les dérives relatives
\[
    \frac{|T(t_k) - T(0)|}{T(0)},
    \qquad
    \frac{\|\vect{L}_P(t_k) - \vect{L}_P(0)\|}{\|\vect{L}_P(0)\|}
\]
qui doivent rester au niveau de l'erreur du schéma, c'est-à-dire l'erreur doit diminuer lorsque l'on raffine le pas de calcule, si elle stagne, cela veut dire qu'il y a une erreur dans le modèle. 

\subsubsection{Non-glissement}

La contrainte $\dot r = -\,\mat{A}_\alpha\,\dot\alpha$ est vérifiée par
construction. Toutefois, il faut que : 
\[
    \|\vect{\lambda}_k\| = \sqrt{\lambda_{1,k}^2 + \lambda_{2,k}^2}
    \leq \mu_s\,m\,g
\]
à chaque pas. Sinon, la sphère glisse et le modèle ne s'applique plus.

\subsubsection{Validité des angles d'Euler}

Le système réduit n'est défini que si $\cos\theta \neq 0$.
\begin{hypotheses}[Validité de la simulation]
    La simulation est valide tant que, à chaque pas,
    \begin{itemize}
        \item $T$ et $\vect{L}_P$ sont conservés à la précision du
        schéma.
        \item $\|\vect{\lambda}_k\| \leq \mu_s\,m\,g$.
        \item $|\cos\theta| > \varepsilon$.
    \end{itemize}
\end{hypotheses}
% ----------------------------------------------------------------------------
\subsection{Architecture de la simulation}

Voici l'architecture de la simulation prête à implémenter avec les rotations et l'orientation représentées par les angles d'Euler.

\usetikzlibrary{fit, backgrounds}
\begin{figure}[H]
    \centering
    \begin{tikzpicture}[>=Stealth,
        bloc/.style={draw, thick, rounded corners=2pt, align=center,
            fill=coulprop!6, draw=coulprop, inner sep=4pt,
            font=\small},
        ctrl/.style={bloc, fill=couldef!8, draw=couldef},
        stop/.style={bloc, fill=coulthm!8, draw=coulthm},
        es/.style={bloc, fill=gray!8, draw=gray},
        fl/.style={->, thick},
        lab/.style={font=\footnotesize, fill=white, inner sep=1pt}]

        \node[es] (in) at (0,0) {Entrées : $m$, $a$, $\mat{I}_b$,
            $\mu_s$ ; $r_0$, $\alpha_0$, $\vect{\Omega}_0$ ; $h$, $N$};
        \node[bloc] (init) at (0,-1.3)
            {Initialisation : $\dot\alpha_0 = \mat{E}_s(\alpha_0)^{-1}
            \vect{\Omega}_0$, \quad $Y_0 = (\alpha_0,\dot\alpha_0)$};
        \node[bloc] (f) at (0,-3.1)
            {Évaluation de $f$ en $Y^{(i)}$, $i = 1,\dots,4$ :\\[2pt]
            $\mat{E}$, $\dot{\mat{E}}\,\dot\alpha$, $\vect{\omega}$
            $\;\to\;$ $\mat{M}_\alpha$, $\vect{h}_\alpha$
            $\;\to\;$ $\mat{A}_\alpha$, $\dot{\mat{A}}_\alpha\,\dot\alpha$
            $\;\to\;$ $\mat{K}_\alpha\,\ddot\alpha^{(i)} = \vect{b}$};
        \node[bloc] (rk) at (-2.9,-4.9)
            {Étape 1 : $Y_{k+1}$ par RK4\\[2pt]
            Étape 2 : $\ddot\alpha_{k+1}$};
        \node[bloc] (rl) at (2.9,-4.9)
            {Étape 3 : $r_{k+1}$ avec les $\vect{g}_i$,\\[2pt]
            $\vect{\lambda}_{k+1}$};
        \node[ctrl] (c) at (0,-6.6)
            {Contrôles : $T$, $\vect{L}_P$ conservés ;
            $\|\vect{\lambda}_{k+1}\| \leq \mu_s\,m\,g$ ;
            $|\cos\theta| > \varepsilon$};
        \node[stop] (arret) at (-3.6,-8.5)
            {Arrêt : simulation\\non valide};
        \node[es] (out) at (3.6,-8.5)
            {Sorties : $r_k$, $\alpha_k$, $\vect{\lambda}_k$,\\
            $k = 0,\dots,N$};

        \draw[fl] (in) -- (init);
        \draw[fl] (init) -- (f);
        \draw[fl] (f) -- (rk);
        \draw[fl] (f) -- (rl);
        \draw[fl] (rk) -- (c);
        \draw[fl] (rl) -- (c);
        \draw[fl, coulthm] (c) -- node[lab, pos=0.45, left=2pt]
            {non} (arret);
        \draw[fl] (c) -- node[lab, pos=0.45, right=2pt]
            {oui, $k = N$} (out);
        \coordinate (droite) at (6.4,0);
        \draw[fl] (c.east) -- (c.east -| droite) -- (f.east -| droite)
            node[lab, pos=0.5, left=2pt, align=right]
            {oui,\\$k \leftarrow k+1$} -- (f.east);

        \begin{scope}[on background layer]
            \node[draw=gray, dashed, rounded corners, inner sep=7pt,
                fit=(f) (rk) (rl) (c) (c.east -| droite)] (boucle) {};
        \end{scope}
        \node[font=\footnotesize, gray, anchor=south west]
            at (boucle.north west) {boucle sur $k$};
    \end{tikzpicture}
    \caption{Architecture de la simulation de la sphère de Chaplygin}
\end{figure}

